% ch30_enriched_day_convolution.tex — Volume III Chapter 6

\chapter{Enriched Categories: Day Convolution}

\begin{overviewbox}
Chapter~20 developed enriched category theory: $\mathcal{V}$-categories,
enriched Yoneda, weighted limits, base change, self-enrichment. Chapter~26
delivered the coend calculus that those constructions depend on.
Chapter~29 identified free monoidal categories with presheaf categories
``with the right tensor.'' This chapter delivers that tensor: \term{Day
convolution}, the canonical monoidal structure on $[\mathcal{C}^{\mathrm{op}},
\mathcal{V}]_\mathcal{V}$ when $\mathcal{C}$ is itself a $\mathcal{V}$-monoidal
category.

Day convolution is a single coend formula that turns the Yoneda embedding
into a strong monoidal functor and exhibits the presheaf category as the
free $\mathcal{V}$-cocompletion-with-monoidal-structure of $\mathcal{C}$.
We develop its construction, prove its universal property, derive the
closed structure, and apply it to enriched Kan extensions and to the
formal theory of monads (preview, Chapter~32). The chapter closes with
Cauchy completion in the Day-convolution-aware sense and a sketch of
where ``monoidal Yoneda'' takes us next.
\end{overviewbox}


% ═══════════════════════════════════════════════════════════════════════════
\section{The Day Convolution Construction}
\label{sec:day-convolution}
% ═══════════════════════════════════════════════════════════════════════════

\subsection{Formalizing}

Fix a complete and cocomplete symmetric monoidal closed category $\mathcal{V}$
and a small $\mathcal{V}$-monoidal category $\mathcal{C}$. That is,
$\mathcal{C}$ is a $\mathcal{V}$-category equipped with a $\mathcal{V}$-functor
$\otimes_\mathcal{C} : \mathcal{C} \otimes \mathcal{C} \to \mathcal{C}$ and a
unit object $I_\mathcal{C} \in \mathcal{C}$, satisfying the $\mathcal{V}$-enriched
monoidal axioms.

\begin{definition}[Day convolution]
\label{def:day-convolution}
For $\mathcal{V}$-functors $F, G : \mathcal{C}^{\mathrm{op}} \to \mathcal{V}$,
the \term{Day convolution} $F \star_{\mathrm{Day}} G : \mathcal{C}^{\mathrm{op}}
\to \mathcal{V}$ is defined by the coend
\begin{equation*}
  (F \star_{\mathrm{Day}} G)(c) \;=\; \int^{a, b \in \mathcal{C}}
  \mathcal{C}(c, a \otimes_\mathcal{C} b) \otimes F(a) \otimes G(b).
\end{equation*}
The Day-convolution unit is the representable $\mathcal{C}(-, I_\mathcal{C})
= Y(I_\mathcal{C})$.
\end{definition}

\begin{howtosay}
$F \star_{\mathrm{Day}} G$ \sayit{``F Day-convolution G''} or simply
\sayit{``F star G''} when the context is clear. The integral signs are over
the pair $(a, b)$ — i.e., a double coend over $\mathcal{C}$ in both arguments.
\end{howtosay}

\subsection{Reorganizing: Why This Formula?}

\begin{remark}[Conceptual reading]
The Day convolution is the ``left Kan extension along the monoidal product''
of the pair $(F, G)$. Explicitly: $\mathcal{C} \otimes \mathcal{C}
\xrightarrow{F \otimes G} \mathcal{V} \otimes \mathcal{V} \xrightarrow{\otimes_\mathcal{V}}
\mathcal{V}$, then take the left Kan extension along $\otimes_\mathcal{C} :
\mathcal{C} \otimes \mathcal{C} \to \mathcal{C}$. The coend formula
(Chapter~26, Example~\ref{ex:lan-coend}) gives the displayed integral.
\end{remark}

\begin{theorem}[Day convolution is a $\mathcal{V}$-functor monoidal structure]
\label{thm:day-monoidal}
$\star_{\mathrm{Day}}$ makes $[\mathcal{C}^{\mathrm{op}}, \mathcal{V}]_\mathcal{V}$
into a $\mathcal{V}$-enriched monoidal category. Its associator and unitors
are induced from those of $\mathcal{C}$ and $\mathcal{V}$.
\end{theorem}

\begin{prooflog}
\proofstep{Show $F \star_{\mathrm{Day}} (G \star_{\mathrm{Day}} H) \cong
(F \star_{\mathrm{Day}} G) \star_{\mathrm{Day}} H$ via the associator of
$\mathcal{C}$.}{Compute both sides by Fubini and the coend manipulations.}
\proofstep{$(F \star_{\mathrm{Day}} (G \star_{\mathrm{Day}} H))(c) =
\int^{a, b} \mathcal{C}(c, a \otimes b) \otimes F(a) \otimes \int^{x, y}
\mathcal{C}(b, x \otimes y) \otimes G(x) \otimes H(y)$.}{Substitute the
definition.}
\proofstep{Apply Fubini (Theorem~\ref{thm:fubini-coend}) and the ninja
Yoneda lemma to integrate over $b$ against $\mathcal{C}(b, x \otimes y)$.}{
Yoneda reduction reduces $\int^b \mathcal{C}(c, a \otimes b) \otimes \mathcal{C}(b, x \otimes y)$
to $\mathcal{C}(c, a \otimes (x \otimes y))$.}
\proofstep{Result: $\int^{a, x, y} \mathcal{C}(c, a \otimes (x \otimes y))
\otimes F(a) \otimes G(x) \otimes H(y)$.}{After Yoneda reduction.}
\proofstep{Symmetrically, $((F \star_{\mathrm{Day}} G) \star_{\mathrm{Day}} H)(c)
\cong \int^{a, x, y} \mathcal{C}(c, (a \otimes x) \otimes y) \otimes F(a)
\otimes G(x) \otimes H(y)$.}{By the same calculation with the inner coend
in the first argument.}
\proofstep{The associator of $\mathcal{C}$ gives an isomorphism
$\mathcal{C}(c, a \otimes (x \otimes y)) \cong \mathcal{C}(c, (a \otimes x)
\otimes y)$, natural in $a, x, y, c$.}{Functoriality of $\mathcal{C}$ along
$\alpha_\mathcal{C}$.}
\proofstep{This induces $F \star_{\mathrm{Day}} (G \star_{\mathrm{Day}} H)
\cong (F \star_{\mathrm{Day}} G) \star_{\mathrm{Day}} H$.}{Coend
functoriality.}
\proofstep{Unit: $Y(I_\mathcal{C}) \star_{\mathrm{Day}} F \cong F$
via the unit-law isomorphism of $\mathcal{C}$. Symmetrically for the right unit.}{
Same Yoneda-reduction argument.}
\proofstep{Pentagon and triangle axioms for $\star_{\mathrm{Day}}$ follow
from those for $\otimes_\mathcal{C}$.}{Functoriality transfers commutativity.}
\end{prooflog}

\subsection{Formally}

\begin{theorem}[Day convolution is closed]
\label{thm:day-closed}
The Day-monoidal category $([\mathcal{C}^{\mathrm{op}}, \mathcal{V}]_\mathcal{V},
\star_{\mathrm{Day}}, Y(I_\mathcal{C}))$ is symmetric monoidal closed when
$\mathcal{C}$ is symmetric monoidal. The internal hom is
\begin{equation*}
  [G, H]_{\mathrm{Day}}(c) \;=\; \int_b \mathcal{V}(G(b), H(c \otimes b)).
\end{equation*}
\end{theorem}

\begin{proof}[Sketch]
Verify the adjunction $[\mathcal{C}^{\mathrm{op}}, \mathcal{V}]_\mathcal{V}
(F \star_{\mathrm{Day}} G, H) \cong [\mathcal{C}^{\mathrm{op}}, \mathcal{V}]_\mathcal{V}
(F, [G, H]_{\mathrm{Day}})$ by direct calculation: substitute the coend formula
for the convolution, use Yoneda reduction on the $\mathcal{C}(c, a \otimes b)$
factor, and identify the resulting double end with $[G, H]_{\mathrm{Day}}$.
Symmetry comes from that of $\mathcal{C}$.
\end{proof}


% ═══════════════════════════════════════════════════════════════════════════
\section{The Universal Property}
\label{sec:day-universal}
% ═══════════════════════════════════════════════════════════════════════════

\subsection{Formalizing}

Day convolution exhibits the presheaf category as the \emph{free
$\mathcal{V}$-cocompletion that is also a $\mathcal{V}$-monoidal extension}
of $\mathcal{C}$.

\begin{theorem}[Universal property of Day convolution]
\label{thm:day-universal}
Let $\mathcal{C}$ be a small $\mathcal{V}$-monoidal category and let
$(\mathcal{E}, \otimes_\mathcal{E}, I_\mathcal{E})$ be a cocomplete
$\mathcal{V}$-monoidal category whose tensor product preserves colimits in
each variable. Then every $\mathcal{V}$-monoidal functor $F : \mathcal{C} \to
\mathcal{E}$ extends, up to essentially unique $\mathcal{V}$-monoidal natural
isomorphism, to a $\mathcal{V}$-cocontinuous strong monoidal functor
\begin{equation*}
  \widetilde F \;:\; ([\mathcal{C}^{\mathrm{op}}, \mathcal{V}]_\mathcal{V},
  \star_{\mathrm{Day}}) \to \mathcal{E}
\end{equation*}
with $\widetilde F \circ Y \cong F$ as $\mathcal{V}$-monoidal functors.
\end{theorem}

\begin{prooflog}
\proofstep{Define $\widetilde F (P) = \int^c P(c) \otimes_\mathcal{E} F(c)$
(weighted colimit; Chapter~28, Theorem~\ref{thm:free-cocompletion}).}{
This is the $\mathcal{V}$-cocontinuous extension.}
\proofstep{Show $\widetilde F$ is strong monoidal: $\widetilde F(F_1 \star_{\mathrm{Day}}
F_2) \cong \widetilde F(F_1) \otimes_\mathcal{E} \widetilde F(F_2)$.}{
The key computation.}
\proofstep{$\widetilde F(F_1 \star_{\mathrm{Day}} F_2) = \int^c (F_1 \star
F_2)(c) \otimes_\mathcal{E} F(c) = \int^{c, a, b} \mathcal{C}(c, a \otimes b)
\otimes_\mathcal{V} F_1(a) \otimes_\mathcal{V} F_2(b) \otimes_\mathcal{E}
F(c)$.}{Substituting definitions.}
\proofstep{Apply Yoneda reduction to the integral over $c$ against
$\mathcal{C}(c, a \otimes b)$: $\int^c \mathcal{C}(c, a \otimes b)
\otimes_\mathcal{E} F(c) \cong F(a \otimes b) \cong F(a) \otimes_\mathcal{E} F(b)$.}{
$F$ is monoidal.}
\proofstep{So $\widetilde F(F_1 \star_{\mathrm{Day}} F_2) \cong \int^{a, b}
F_1(a) \otimes_\mathcal{V} F_2(b) \otimes_\mathcal{E} F(a) \otimes_\mathcal{E} F(b)$.}{
Substituting and rearranging.}
\proofstep{Right-hand side: $\widetilde F(F_1) \otimes_\mathcal{E}
\widetilde F(F_2) = \left(\int^a F_1(a) \otimes_\mathcal{E} F(a)\right)
\otimes_\mathcal{E} \left(\int^b F_2(b) \otimes_\mathcal{E} F(b)\right)$.}{
Substituting the definition of $\widetilde F$.}
\proofstep{Since $\otimes_\mathcal{E}$ preserves colimits, we can move the
$\otimes_\mathcal{E}$ inside the coends; the two expressions agree.}{
Cocontinuity of $\otimes_\mathcal{E}$ plus Fubini.}
\proofstep{Unit: $\widetilde F(Y(I_\mathcal{C})) = \int^c \mathcal{C}(c, I_\mathcal{C})
\otimes_\mathcal{E} F(c) \cong F(I_\mathcal{C}) \cong I_\mathcal{E}$.}{Yoneda
reduction then $F$ preserves the unit.}
\proofstep{Uniqueness: any cocontinuous strong monoidal $G$ extending $F$
must agree with $\widetilde F$ on representables (by $G \circ Y \cong F$),
hence on all presheaves (by cocontinuity and density).}{Standard uniqueness
argument.}
\end{prooflog}

\subsection{Reorganizing: The Slogan}

\begin{corollary}[Day convolution = free monoidal cocompletion]
\label{cor:day-free}
For a small $\mathcal{V}$-monoidal category $\mathcal{C}$, the Yoneda
embedding $Y : \mathcal{C} \to [\mathcal{C}^{\mathrm{op}}, \mathcal{V}]_\mathcal{V}$
exhibits the Day-monoidal presheaf category as the free $\mathcal{V}$-cocomplete
monoidal extension of $\mathcal{C}$.
\end{corollary}

\begin{remark}[The full picture]
Combining Chapter~28 (Yoneda as free cocompletion) and the present theorem:
\begin{itemize}
  \item Cocompletion alone: $\mathcal{C} \xrightarrow{Y} [\mathcal{C}^{\mathrm{op}},
        \mathcal{V}]_\mathcal{V}$. Universal among cocomplete extensions.
  \item Monoidal cocompletion: same target with Day convolution. Universal
        among cocomplete monoidal extensions.
\end{itemize}
Day convolution is the canonical lift of the universal property from
``cocompletion'' to ``monoidal cocompletion.''
\end{remark}


% ═══════════════════════════════════════════════════════════════════════════
\section{Enriched Kan Extensions}
\label{sec:enriched-kan-via-day}
% ═══════════════════════════════════════════════════════════════════════════

The coend calculus and Day convolution combine to give the cleanest
formulation of enriched Kan extensions.

\subsection{Formalizing}

\begin{theorem}[Enriched Kan extension formulas]
\label{thm:enriched-kan}
For $\mathcal{V}$-functors $K : \mathcal{C} \to \mathcal{D}$ and $F :
\mathcal{C} \to \mathcal{E}$ (with $\mathcal{E}$ tensored and cocomplete),
the \term{enriched left Kan extension} $\operatorname{Lan}_K F : \mathcal{D}
\to \mathcal{E}$ exists and is given by
\begin{equation*}
  (\operatorname{Lan}_K F)(d) \;=\; \int^c \mathcal{D}(K c, d) \otimes_\mathcal{E} F(c).
\end{equation*}
Dually, the \term{enriched right Kan extension} $\operatorname{Ran}_K F$
(with $\mathcal{E}$ cotensored and complete) is
\begin{equation*}
  (\operatorname{Ran}_K F)(d) \;=\; \int_c [\mathcal{D}(d, K c), F(c)]_\mathcal{E}.
\end{equation*}
\end{theorem}

\begin{proof}
The formulas are the weighted colimit and weighted limit, respectively
(Chapter~27, \S~\ref{sec:wkan-preview}). Existence follows from cocompleteness
(resp.\ completeness) of $\mathcal{E}$ via the reduction of weighted
(co)limits to conical (co)limits + (co)tensors (Chapter~27, Theorem~\ref{thm:wcomplete-from-pieces}).
\end{proof}

\begin{theorem}[Kan extensions form an adjoint triple along Yoneda]
\label{thm:kan-triple-along-yoneda}
For the Yoneda embedding $Y : \mathcal{C} \to [\mathcal{C}^{\mathrm{op}},
\mathcal{V}]_\mathcal{V}$, there is an adjoint triple
\begin{equation*}
  \operatorname{Lan}_Y \;\dashv\; Y^* \;\dashv\; \operatorname{Ran}_Y
\end{equation*}
between $[\mathcal{C}, \mathcal{E}]_\mathcal{V}$ (functor $\mathcal{V}$-category)
and $[[\mathcal{C}^{\mathrm{op}}, \mathcal{V}]_\mathcal{V}, \mathcal{E}]_\mathcal{V}$
(presheaf functor category). Here $Y^*$ is precomposition with $Y$.
\end{theorem}

\begin{proof}[Sketch]
This is the enriched analogue of Chapter~18's adjoint triple
$\operatorname{Lan}_p \dashv p^* \dashv \operatorname{Ran}_p$ specialized to
$p = Y$. The proof uses the coend formulas of Theorem~\ref{thm:enriched-kan}
and the universal property of the presheaf category.
\end{proof}

\subsection{Reorganizing: Day Convolution as Kan Extension}

\begin{proposition}[Day convolution as $\operatorname{Lan}$]
\label{prop:day-as-lan}
The Day convolution $F \star_{\mathrm{Day}} G$ is the left Kan extension of
the composite
\begin{equation*}
  \mathcal{C} \otimes \mathcal{C} \xrightarrow{Y \otimes Y}
  [\mathcal{C}^{\mathrm{op}}, \mathcal{V}]_\mathcal{V} \otimes [\mathcal{C}^{\mathrm{op}},
  \mathcal{V}]_\mathcal{V} \xrightarrow{\star_{\mathrm{Day}}}
  [\mathcal{C}^{\mathrm{op}}, \mathcal{V}]_\mathcal{V}
\end{equation*}
along $\otimes_\mathcal{C} : \mathcal{C} \otimes \mathcal{C} \to \mathcal{C}$,
evaluated at the pair $(F, G)$.
\end{proposition}

\begin{proof}
Direct from Definition~\ref{def:day-convolution} and the left Kan coend
formula of Theorem~\ref{thm:enriched-kan}.
\end{proof}


% ═══════════════════════════════════════════════════════════════════════════
\section{Monoids in the Day-Monoidal Category}
\label{sec:day-monoids}
% ═══════════════════════════════════════════════════════════════════════════

The most striking applications of Day convolution come from looking at
\emph{monoids} in the resulting monoidal category $[\mathcal{C}^{\mathrm{op}},
\mathcal{V}]_\mathcal{V}$.

\subsection{Formalizing}

\begin{theorem}[Monoids in Day convolution]
\label{thm:day-monoids}
For a small $\mathcal{V}$-monoidal category $\mathcal{C}$:
\begin{equation*}
  \mathrm{Mon}([\mathcal{C}^{\mathrm{op}}, \mathcal{V}]_\mathcal{V},
  \star_{\mathrm{Day}})
  \;\simeq\;
  [\mathcal{C}, \mathcal{V}]_\mathcal{V}^{\mathrm{lax\,monoidal}},
\end{equation*}
the category of \term{lax monoidal $\mathcal{V}$-functors} $\mathcal{C} \to
\mathcal{V}$.
\end{theorem}

\begin{remark}[Interpretation]
A monoid in $[\mathcal{C}^{\mathrm{op}}, \mathcal{V}]_\mathcal{V}$ with Day
tensor is precisely a lax monoidal functor $\mathcal{C} \to \mathcal{V}$.
This is the \emph{universal monoid construction}: Day convolution turns
``lax monoidal functor'' into ``monoid in a single category.''
\end{remark}

\subsection{Reorganizing: Examples and Special Cases}

\begin{example_thm}[Algebras over a $\mathcal{V}$-operad]
For an operad $\mathcal{O}$ in $\mathcal{V}$, monoids in the Day-monoidal
$[\mathbb{F}^{\mathrm{op}}, \mathcal{V}]$ (with $\mathbb{F}$ the category of
finite sets and bijections, monoidal under disjoint union) recover
$\mathcal{O}$-algebras. Operads-as-monoids becomes a special case of
Day-convolution monoid theory.
\end{example_thm}

\begin{example_thm}[The Hochschild cohomology setup]
For an associative algebra $A$, the Hochschild cohomology of $A$ arises as
the cohomology of the Day-monoidal presheaf category $[\Delta^{\mathrm{op}},
\mathcal{V}]$ at the specific monoid $A$. Day convolution captures the
``cosimplicial homotopical structure'' uniformly.
\end{example_thm}


% ═══════════════════════════════════════════════════════════════════════════
\section{Cauchy Completion Revisited}
\label{sec:cauchy-revisited}
% ═══════════════════════════════════════════════════════════════════════════

Chapter~28 introduced Cauchy completeness in the (unenriched) profunctor
setting. In the enriched setting, Cauchy completeness interacts richly with
Day convolution.

\subsection{Formalizing}

\begin{definition}[Enriched Cauchy completion]
\label{def:enriched-cauchy}
A $\mathcal{V}$-category $\mathcal{C}$ is \term{Cauchy complete} if for
every $\mathcal{V}$-profunctor $P : \mathbf{1} \to \mathcal{C}$ admitting a
right adjoint as a 1-cell in $\mathcal{V}$-$\mathbf{Prof}$, $P$ is
representable by an object of $\mathcal{C}$.

The \term{Cauchy completion} $\bar{\mathcal{C}}$ is the universal Cauchy-complete
enlargement; it is precisely the full sub-$\mathcal{V}$-category of
$[\mathcal{C}^{\mathrm{op}}, \mathcal{V}]_\mathcal{V}$ consisting of the
\term{absolute presheaves}.
\end{definition}

\begin{remark}[Concrete cases]
\begin{itemize}
  \item In the $\mathbf{Set}$-enriched case, Cauchy completion is the
        idempotent (Karoubian) completion.
  \item In the $\mathbf{Ab}$-enriched case, Cauchy completion adds direct
        summands of free objects.
  \item In the $\mathbf{Cat}$-enriched case (i.e., 2-categories), Cauchy
        completion adds objects represented by 1-cells with a right adjoint.
\end{itemize}
\end{remark}


% ═══════════════════════════════════════════════════════════════════════════
\section{Exercises}
\label{sec:day-exercises}
% ═══════════════════════════════════════════════════════════════════════════

\begin{exercisesbox}
\begin{qa}
\qaitem{Define Day convolution.}{For $F, G : \mathcal{C}^{\mathrm{op}} \to \mathcal{V}$, $(F \star_{\mathrm{Day}} G)(c) = \int^{a, b} \mathcal{C}(c, a \otimes_\mathcal{C} b) \otimes F(a) \otimes G(b)$.}
\qaitem{What is the Day-convolution unit?}{The Yoneda embedding of the monoidal unit: $Y(I_\mathcal{C}) = \mathcal{C}(-, I_\mathcal{C})$.}
\qaitem{Why is Day convolution monoidal?}{Because the associator and unitors of $\mathcal{C}$ transfer (via Yoneda reduction and functoriality of coends) to associator and unitors of $\star_{\mathrm{Day}}$.}
\qaitem{Show that $Y(I_\mathcal{C}) \star_{\mathrm{Day}} F \cong F$.}{$(Y(I_\mathcal{C}) \star F)(c) = \int^{a, b} \mathcal{C}(c, a \otimes b) \otimes \mathcal{C}(a, I) \otimes F(b)$; Yoneda-reduce $a$ to get $\int^b \mathcal{C}(c, I \otimes b) \otimes F(b) \cong \int^b \mathcal{C}(c, b) \otimes F(b) \cong F(c)$.}
\qaitem{What is the internal hom of Day convolution?}{$[G, H]_{\mathrm{Day}}(c) = \int_b \mathcal{V}(G(b), H(c \otimes b))$.}
\qaitem{When is Day convolution closed?}{When $\mathcal{C}$ is symmetric monoidal and $\mathcal{V}$ is symmetric monoidal closed and complete.}
\qaitem{State the universal property of Day convolution.}{For every cocomplete $\mathcal{V}$-monoidal $\mathcal{E}$ and $\mathcal{V}$-monoidal $F : \mathcal{C} \to \mathcal{E}$, there is an essentially unique cocontinuous strong monoidal $\widetilde F : [\mathcal{C}^{\mathrm{op}}, \mathcal{V}]_\mathcal{V} \to \mathcal{E}$ extending $F$ via Yoneda.}
\qaitem{Why is Day convolution the ``free monoidal cocompletion''?}{Because $[\mathcal{C}^{\mathrm{op}}, \mathcal{V}]_\mathcal{V}$ with Day tensor is initial among cocomplete monoidal extensions of $\mathcal{C}$; equivalently, the universal property uniquely characterizes the construction.}
\qaitem{Write the enriched left Kan extension formula.}{$(\operatorname{Lan}_K F)(d) = \int^c \mathcal{D}(K c, d) \otimes_\mathcal{E} F(c)$.}
\qaitem{Right Kan extension formula?}{$(\operatorname{Ran}_K F)(d) = \int_c [\mathcal{D}(d, K c), F(c)]_\mathcal{E}$.}
\qaitem{Why is $\operatorname{Lan}_Y \dashv Y^* \dashv \operatorname{Ran}_Y$ an adjoint triple?}{Because Yoneda is the universal embedding into the free cocompletion; the restriction $Y^*$ has left and right adjoints by the colimit/limit formulas (special case of Chapter~18's adjoint triple).}
\qaitem{Why is Day convolution the left Kan extension of $\star_\mathcal{V} \circ (Y \otimes Y)$ along $\otimes_\mathcal{C}$?}{Because Yoneda extends pointwise to the presheaf category, and the Day formula is exactly the left Kan extension of the composite ``embed-then-tensor'' along $\otimes_\mathcal{C}$.}
\qaitem{State the monoids-in-Day-convolution theorem.}{Monoids in $([\mathcal{C}^{\mathrm{op}}, \mathcal{V}]_\mathcal{V}, \star_{\mathrm{Day}})$ are equivalent to lax monoidal $\mathcal{V}$-functors $\mathcal{C} \to \mathcal{V}$.}
\qaitem{Interpret this for operads.}{For $\mathcal{C} = \mathbb{F}$ (finite sets and bijections, disjoint union), monoids in Day convolution = lax monoidal $\mathbb{F} \to \mathcal{V}$ = operads.}
\qaitem{Why is Day convolution natural in the base category?}{Because the formula's coend integrates over all morphisms; functoriality in $\mathcal{C}$ is preserved by the coend construction.}
\qaitem{What is an enriched profunctor $\mathcal{C} \to \mathcal{D}$?}{A $\mathcal{V}$-functor $P : \mathcal{D}^{\mathrm{op}} \otimes \mathcal{C} \to \mathcal{V}$.}
\qaitem{What is enriched Cauchy completion?}{The universal Cauchy-complete enlargement: full sub-$\mathcal{V}$-category of presheaves consisting of absolute (representable-up-to-equivalence) ones.}
\qaitem{What is $\mathbf{Ab}$-Cauchy completion?}{Adding direct summands of free objects (i.e., the idempotent splitting for direct sum decompositions).}
\qaitem{What does ``absolute'' mean for a presheaf in this context?}{A presheaf $P$ is absolute if its Yoneda extension $\widetilde Y^{-1}(P)$ exists in $\mathcal{C}$ — i.e., $P$ is represented by an object of the Cauchy completion.}
\qaitem{Why is the Day-convolution monoidal structure ``the right one''?}{Because it is uniquely determined by the universal property: it is the unique monoidal structure on the free cocompletion that makes Yoneda strongly monoidal.}
\qaitem{What is the relationship between Day convolution and the convolution of $L^1$-functions on a group?}{Same formula: $(f * g)(x) = \int f(a) g(x - a) \, da$ is the special case for $\mathcal{C}$ = a group and $\mathcal{V} = \mathbf{Set}$; the coend becomes an ordinary integral over the group.}
\qaitem{State the Day convolution formula for a discrete (groupoid) $\mathcal{C}$.}{$(F \star_{\mathrm{Day}} G)(c) = \coprod_{a \otimes b \cong c} F(a) \times G(b) / \sim$, where the coend reduces to a coproduct over a single representative; in a group, this is the convolution sum.}
\qaitem{Why does $\otimes_\mathcal{E}$ need to preserve colimits for the universal property to hold?}{Because $\widetilde F$ is a colimit-preserving functor; for it to be \emph{monoidal}, the target tensor must distribute over colimits. Otherwise the convolution formula doesn't carry the monoidal structure cleanly.}
\qaitem{What's the simplest example of Day convolution?}{$\mathcal{C} = $ a discrete monoidal category (a monoid). Then $[\mathcal{C}^{\mathrm{op}}, \mathbf{Set}]$ is the category of $\mathcal{C}$-graded sets, and Day convolution is exactly the convolution on graded sets.}
\qaitem{What is the dual of the Day convolution theorem (using right Kan extension)?}{The cofree comonoidal cocompletion: $[\mathcal{C}, \mathcal{V}]_\mathcal{V}$ with the dual ``Day cotensor'' has the universal property for comonoids and oplax monoidal functors.}
\qaitem{How does Day convolution relate to Chapter~29's free monoidal category?}{$F^\otimes(\mathcal{C}) \cong ([\mathcal{C}^{\mathrm{op}}, \mathbf{Set}], \star_{\mathrm{Day}})$ when $\mathcal{C}$ is itself monoidal; Day convolution provides the missing tensor on the presheaf category.}
\qaitem{What is the explicit formula for an $\mathcal{O}$-algebra in $\mathcal{V}$ via Day convolution?}{An $\mathcal{O}$-algebra $A$ corresponds to a lax monoidal functor $\mathbb{F} \to \mathcal{V}$ sending $n$ to $A^{\otimes n}$, with the operad's composition giving the structure 2-cells.}
\qaitem{State Day convolution closure as an adjunction.}{$[\mathcal{C}^{\mathrm{op}}, \mathcal{V}]_\mathcal{V}(F \star_{\mathrm{Day}} G, H) \cong [\mathcal{C}^{\mathrm{op}}, \mathcal{V}]_\mathcal{V}(F, [G, H]_{\mathrm{Day}})$.}
\qaitem{Why is the proof of associativity for Day convolution a coend calculation?}{Because the formula is itself a coend; associativity reduces to Fubini-plus-Yoneda-reduction, the two pillars of the coend calculus.}
\qaitem{What chapter develops the formal theory of monads that uses Day convolution?}{Chapter~32 (The Formal Theory of Monads); monoids in Day-monoidal categories provide the foundation for monad transformers and distributive laws (Chapter~35).}
\end{qa}
\end{exercisesbox}


% ═══════════════════════════════════════════════════════════════════════════
\section*{Chapter Summary}
\addcontentsline{toc}{section}{Chapter Summary}
% ═══════════════════════════════════════════════════════════════════════════

\begin{summarybox}
\textbf{Day convolution.}\quad For a $\mathcal{V}$-monoidal $\mathcal{C}$,
the convolution $(F \star_{\mathrm{Day}} G)(c) = \int^{a, b} \mathcal{C}(c,
a \otimes b) \otimes F(a) \otimes G(b)$ makes $[\mathcal{C}^{\mathrm{op}},
\mathcal{V}]_\mathcal{V}$ into a $\mathcal{V}$-monoidal category. The unit is
$Y(I_\mathcal{C})$.

\textbf{Universal property.}\quad The Yoneda embedding $Y : \mathcal{C} \to
[\mathcal{C}^{\mathrm{op}}, \mathcal{V}]_\mathcal{V}$ exhibits the
Day-monoidal presheaf category as the free $\mathcal{V}$-cocomplete monoidal
extension of $\mathcal{C}$: any $\mathcal{V}$-monoidal $F : \mathcal{C} \to
\mathcal{E}$ extends uniquely to a cocontinuous strong monoidal $\widetilde F$.

\textbf{Closed structure.}\quad For symmetric monoidal $\mathcal{C}$, Day
convolution is symmetric monoidal closed, with internal hom $[G, H]_{\mathrm{Day}}(c)
= \int_b \mathcal{V}(G(b), H(c \otimes b))$.

\textbf{Enriched Kan extensions.}\quad $\operatorname{Lan}_K F (d) = \int^c
\mathcal{D}(K c, d) \otimes F(c)$ and $\operatorname{Ran}_K F (d) = \int_c
[\mathcal{D}(d, K c), F(c)]$. These give the adjoint triple
$\operatorname{Lan}_Y \dashv Y^* \dashv \operatorname{Ran}_Y$ along the
Yoneda embedding.

\textbf{Day-monoid theorem.}\quad Monoids in $([\mathcal{C}^{\mathrm{op}},
\mathcal{V}]_\mathcal{V}, \star_{\mathrm{Day}})$ are equivalent to lax
monoidal $\mathcal{V}$-functors $\mathcal{C} \to \mathcal{V}$. Operads and
algebraic structures are special cases.

\textbf{Cauchy completion.}\quad The full sub-presheaf category of
\emph{absolute} presheaves. In $\mathbf{Set}$-enriched: idempotent splitting;
in $\mathbf{Ab}$-enriched: direct summand completion; in $\mathbf{Cat}$-enriched:
left-adjoint completion.
\end{summarybox}


% ═══════════════════════════════════════════════════════════════════════════
\section*{Formal Restatement}
\addcontentsline{toc}{section}{Formal Restatement}
% ═══════════════════════════════════════════════════════════════════════════

\begin{formalrestatement}
\textbf{Day convolution.}\quad
\begin{equation*}
  (F \star_{\mathrm{Day}} G)(c) \;=\; \int^{a, b \in \mathcal{C}}
  \mathcal{C}(c, a \otimes_\mathcal{C} b) \otimes F(a) \otimes G(b),
  \qquad \mathbf{1}_{\mathrm{Day}} \;=\; Y(I_\mathcal{C}).
\end{equation*}

\medskip
\textbf{Internal hom (Day-closed).}\quad
\begin{equation*}
  [G, H]_{\mathrm{Day}}(c) \;=\; \int_b \mathcal{V}(G(b), H(c \otimes b)).
\end{equation*}

\medskip
\textbf{Universal property.}\quad The Day-monoidal $[\mathcal{C}^{\mathrm{op}},
\mathcal{V}]_\mathcal{V}$ is the free $\mathcal{V}$-cocomplete monoidal
extension of $\mathcal{C}$. $\widetilde F(P) = \int^c P(c) \otimes_\mathcal{E} F(c)$.

\medskip
\textbf{Enriched Kan extensions.}\quad
\begin{equation*}
  (\operatorname{Lan}_K F)(d) \;=\; \int^c \mathcal{D}(K c, d) \otimes_\mathcal{E} F(c),
  \qquad
  (\operatorname{Ran}_K F)(d) \;=\; \int_c [\mathcal{D}(d, K c), F(c)]_\mathcal{E}.
\end{equation*}
Adjoint triple along $Y$: $\operatorname{Lan}_Y \dashv Y^* \dashv \operatorname{Ran}_Y$.

\medskip
\textbf{Day-monoid theorem.}\quad
\begin{equation*}
  \mathrm{Mon}([\mathcal{C}^{\mathrm{op}}, \mathcal{V}]_\mathcal{V},
  \star_{\mathrm{Day}}) \;\simeq\; [\mathcal{C}, \mathcal{V}]_\mathcal{V}^{\mathrm{lax\,mon}}.
\end{equation*}

\medskip
\textbf{Cauchy completion.}\quad $\bar{\mathcal{C}} \subset [\mathcal{C}^{\mathrm{op}},
\mathcal{V}]_\mathcal{V}$ is the full sub-$\mathcal{V}$-category of absolute
presheaves. Examples: idempotent splitting ($\mathbf{Set}$), direct-summand
closure ($\mathbf{Ab}$), left-adjoint splitting ($\mathbf{Cat}$).

\medskip
\noindent\textit{Chapter~31 develops the general theory of Kan extensions
in the enriched setting, building on the formulas of this chapter and the
adjoint triple along Yoneda. Pointwise extensions, absolute extensions, and
applications to coherence and to the formal theory of monads (Chapter~32)
follow.}
\end{formalrestatement}
